\section{What broke}
\label{sec:whatbroke}

The project keeps a running log of surprises, \code{docs/WHAT\_BROKE.md}, written
as they happened. This section keeps the entries a reader needs to judge the
evaluation. They fall into three groups: ways the evaluation could have leaked
or flattered a method, bugs in rivals and scorers, and data-collection failures
that changed what the pools contain. Each entry says what happened and what was
done.

\subsection{Leaks and flattering rules}

\paragraph{Route ids sort by category.}
Fail2Drive's raw route ids sort by scenario type. In the first design of the
caption-then-cluster hypothesis (H25), two of six grouping runs would have
listed the route descriptions in category blocks, and the answer mapping sat in
the parent folder of the frames. A review caught it before the lock
(\code{WHAT\_BROKE.md}, 2026-11-12). Since then every cache, frame folder and
caption request is ordered by a salted hash, and an audit lists every file each
captioning agent touched. Opaque ids are not blind if their order is not.

\paragraph{A route field gave the oracle a perfect score.}
On the first dev run, B2 with labels recovered PedestriansOnRoad perfectly by
reading a route field that names the type (\code{VARIANTS.md}, v1). B2 has been
an oracle ceiling only since then, never a rival.

\paragraph{Choosing the best axis with the key.}
The early dev rule scored each method by its best axis for each gap, chosen with
the answer key. That flatters every method, and most of all methods that emit
many axes. A1 replaced it with the method's own top 10.

\paragraph{Ward cuts listed the same group twice.}
Cutting a tree at 2, 3 and 5 clusters lists a cluster again whenever a later cut
leaves it whole. Under A1 most methods therefore had 7 or 8 distinct axes in
their top 10, and the number depended on the shape of their tree. Jerry's agent
found it; A2 counts distinct axes only and keeps cutting until there are 10.

\paragraph{A probe can read the window's position.}
m19 chose layers by how well a pass/fail probe separates routes, and the probes
reached an area under the curve near 1.
% measured: AUROC about 0.99 (VARIANTS.md, batch 3, m19)
A passing route's window ends at the route's end and a failing route's at the
failure, so the probe was reading where the window sits, not anything about the
failure. The layer choice was uninformative.

\paragraph{One gap dominated the first dev corpus.}
On the first dev evaluation one practice gap made up most of the failures,
% measured: 51 of 73 failures (WHAT_BROKE.md, 2026-10-15)
51 of 73. An axis holding every failure then scores well, and every method,
random included, scored the same. The evaluation needs pools where each gap is
a minority of failures, which Fail2Drive's pairs give.

\subsection{Bugs in rivals and scorers}

\paragraph{B3-behaviour's window ended early.}
A1 says B3-behaviour describes speed just before the failure. The cache kept
ego speed for the first 100 frames only, while activations used every frame.
% measured: 645 of 1,185 routes; median about 10 s early (WHAT_BROKE.md, 2026-11-16)
On 645 of SimLingo's 1,185 scored routes the rival's window therefore ended a
median of about 10\,s before the true one. The same entry found the window is
about 1.5\,s, not the 3\,s A1 states. The cap was removed before any held-out
score (A3), which can only raise the rival. The 6-frame window stands because
it is what A1 tagged in code.

\paragraph{The chance column drew from passing routes.}
B1 is defined as random axes drawn from failing routes. The sprint scorer's
first chance column drew from every route, so it read too low
(\code{WHAT\_BROKE.md}, 2026-11-15). Rival comparisons and $p$ values never used
that column. The scorer now draws from failing routes only.

\paragraph{The H21b pair unit carried the seed.}
\label{sec:h21b}
When the three H21 seeds were pooled, the pool builder set each route's
\code{pair} to its \code{base\_id}, which carries the seed (for example
\code{f2d-0}, \code{f2d-0-s25}, \code{f2d-0-s26}). The pair bootstrap therefore
resampled pair-seed units instead of Fail2Drive pairs carrying all their seeds,
against A1.
% measured: 299 seed-tagged units vs 100 pairs (results/h21/h21b/manifest.json and routes.jsonl, -sNN stripped)
The 593-route pool has 299 such units against 100 real pairs. Treating copies of
one pair on different seeds as independent made the intervals and $p$ values too
narrow. An independent scorer written from the pre-registration text reproduced
every point score exactly. With the suffix stripped it gave wider intervals;
with the seed-tagged unit it reproduced the logged intervals, so the unit
explains the whole difference. The reading of that replication did not change,
and the UniAD manifest's \code{pair} field has no seed suffix. The correction is
recorded in \code{VARIANTS.md} (H21b replication, ``Correction''); it has no
entry in \code{WHAT\_BROKE.md} yet. Any pool built by concatenating seeds must
strip the suffix before resampling.

\paragraph{Unit tests at the wrong shape.}
Two j methods built on the passing routes' covariance passed their unit tests,
which used fewer columns than passing routes. The real cache has more columns
than passing routes, so that covariance is singular, and one method could only
see failure directions inside the passing routes' span
(\code{WHAT\_BROKE.md}, 2026-11-15). Both now reduce to 32 principal components
first. A test for a covariance method must use the real ratio of columns to rows.

\paragraph{A minimum group size larger than the gaps.}
The first frozen activation method kept only groups of at least 20 failing
routes. H21's categories had
% measured: 8, 5 and 3 failing routes (WHAT_BROKE.md, 2026-11-12)
8, 5 and 3 failures, so it could not score well by construction. This was
knowable from the code and the counts before anything ran.

\subsection{Data collection}

\paragraph{A third of SimLingo's test pool was lost on the first pass.}
% measured: 803 of 1,200 routes with an outcome after the first arrays (WHAT_BROKE.md, 2026-11-09)
Only 803 of 1,200 routes had an outcome after the first arrays. Part was a
cluster-wide slowdown in model loading. Part was ours: a process-cleanup
function returned early, so orphaned processes kept GPU memory and the next
route in the same slot ran out of memory. The cleanup was fixed and tested, and
the lost routes were rerun once, as the pre-registration allows.

\paragraph{TransFuser++'s hooks never fired.}
The TransFuser++ agent calls its network's \code{forward} method directly, which
skips the hooks PyTorch runs on a normal call, so the first capture recorded no
model step (\code{WHAT\_BROKE.md}, 2026-11-05). The adapter now wraps
\code{forward} itself.

\paragraph{Route markers in earlier SimLingo runs.}
Earlier development runs passed a debug flag that draws the route's waypoints and
a trigger point labelled with the scenario's name into the world, visible in the
policy's camera (\code{WHAT\_BROKE.md}, 2026-11-03). Every Fail2Drive test pool
was driven with the flag off, so the pools in this paper carry no markers. The
views pre-registration still checks sampled frames by eye before captioning.

\paragraph{A rerun is a new sample, not a replay.}
Repeated CARLA runs of one route and seed differ from the first frame
(\code{WHAT\_BROKE.md}, 2026-10-06). So the TransFuser++ re-drive that saved
camera frames is a new drive, and the caption row on TransFuser++ covers only
pairs whose failures reproduce, with that coverage stated.

\paragraph{Frames were thought missing.}
The sprint's first checks found no camera frames in either test pool's shards,
so A1 dropped caption methods. On 2026-11-18 it was found that every SimLingo
test route's run folder had saved the model's own camera input every 10 steps
(\code{VARIANTS.md}, checks before the loop, correction). That is why the caption
row could return in the views pre-registration.

\paragraph{UniAD drives far more slowly than planned.}
% measured: median 34 min per driven route over 624 routes (WHAT_BROKE.md, 2026-11-18)
UniAD routes took a median of 34 minutes each, so most array tasks hit their time
limit and only part of each seed was driven on the first pass. The missing
routes were resubmitted. Because the full pool lands late, the views
pre-registration fixed an interim read on the routes driven by a set time, always
labelled interim, beside the final read (C43).
